% Параметры задания: ФИО ученика, Путь А, N=5. :contentReference[oaicite:0]{index=0}

\documentclass[a4paper,12pt]{article}

\usepackage[T2A]{fontenc}
\usepackage[utf8]{inputenc}
\usepackage[russian]{babel}

\usepackage{amsmath,amssymb,mathtools}
\usepackage{helvet}
\renewcommand{\familydefault}{\sfdefault}

\usepackage{icomma}
\usepackage[
    a4paper,
    margin=18mm,
    top=18mm,
    bottom=20mm,
    headheight=25pt
]{geometry}

\usepackage{microtype}
\usepackage{tikz}
\usetikzlibrary{
    calc,
    arrows.meta,
    positioning,
    shapes.geometric
}

\usepackage{pgfplots}
\pgfplotsset{compat=1.18}

\usepackage{tabularx,booktabs}
\usepackage{enumitem}
\usepackage{hyperref}
\usepackage{parskip}
\usepackage{fancyhdr}
\usepackage{setspace}
\usepackage{xcolor}

% ------------------------------------------------
% Цветовая палитра
% ------------------------------------------------

\definecolor{accent}{HTML}{2F6F6D}
\definecolor{darkaccent}{HTML}{214D4B}
\definecolor{textgray}{HTML}{333333}
\definecolor{softgray}{HTML}{E8ECEC}

% ------------------------------------------------
% Общая типографика
% ------------------------------------------------

\hypersetup{
    colorlinks=true,
    linkcolor=darkaccent,
    urlcolor=darkaccent
}

\onehalfspacing
\setlength{\parindent}{0pt}

\setlist[itemize]{
    leftmargin=6mm,
    itemsep=2pt,
    topsep=3pt
}

\setlist[enumerate]{
    leftmargin=7mm,
    itemsep=4pt,
    topsep=3pt
}

% ------------------------------------------------
% Колонтитулы
% ------------------------------------------------

\pagestyle{fancy}
\fancyhf{}

\fancyhead[L]{\small\color{textgray}Линейная функция}
\fancyhead[R]{\small\color{textgray}Ярослав Верещагин}
\fancyfoot[C]{\small\thepage}

\renewcommand{\headrulewidth}{0.3pt}
\renewcommand{\headrule}{%
    \hbox to\headwidth{%
        \color{softgray}%
        \leaders\hrule height \headrulewidth\hfill
    }%
}

% ------------------------------------------------
% Заголовки разделов
% ------------------------------------------------

\newcommand{\topic}[1]{%
    \vspace{0.5em}%
    {\Large\bfseries\color{darkaccent}#1}%
    \par\vspace{0.35em}%
    {\color{softgray}\rule{\linewidth}{0.7pt}}%
    \vspace{0.45em}%
}

% ------------------------------------------------
% Стили блок-схемы
% ------------------------------------------------

\tikzset{
    flowstart/.style={
        ellipse,
        draw=accent,
        very thick,
        minimum width=34mm,
        minimum height=10mm,
        align=center,
        font=\small
    },
    flowaction/.style={
        rounded corners=3pt,
        draw=accent,
        very thick,
        minimum width=72mm,
        minimum height=12mm,
        text width=66mm,
        align=center,
        font=\small
    },
    flowdecision/.style={
        diamond,
        aspect=2.4,
        draw=accent,
        very thick,
        minimum width=45mm,
        minimum height=17mm,
        text width=34mm,
        align=center,
        font=\small
    },
    flowarrow/.style={
        -{Latex[length=3mm]},
        very thick,
        draw=darkaccent
    }
}

\begin{document}

% ================================================================
% ТИТУЛЬНЫЙ ЛИСТ
% ================================================================

\begin{titlepage}

\thispagestyle{empty}
\centering

\vspace*{26mm}

{\Large\color{darkaccent}\bfseries
Чек-лист
\par}

\vspace{8mm}

{\Huge\bfseries
Линейная функция и график прямой
\par}

\vspace{4mm}

{\Large
$y=kx+b$
\par}

\vspace{18mm}

{\large Коэффициенты $k$ и $b$\par}

\vspace{2mm}

{\large Определение формулы по графику\par}

\vspace{2mm}

{\large Семейства $y=m$ и $y=kx$\par}

\vspace{2mm}

{\large Применение к графикам движения\par}

\vfill

{\large
Лёвин Артём Александрович
\par}

\vspace{2mm}

{\large
эксклюзивно для Ярослава Верещагина
\par}

\vspace{5mm}

{\large
\today
\par}

\vspace{10mm}

\begin{tikzpicture}[remember picture,overlay]
    \draw[
        accent!55,
        line width=1.2pt
    ]
    ([xshift=18mm,yshift=18mm]current page.south west)
    ..
    controls
    ([xshift=70mm,yshift=31mm]current page.south west)
    and
    ([xshift=-72mm,yshift=8mm]current page.south east)
    ..
    ([xshift=-18mm,yshift=18mm]current page.south east);
\end{tikzpicture}

\end{titlepage}

% ================================================================
% 1
% ================================================================

\topic{1. Общий вид линейной функции}

Линейная функция задаётся формулой
\[
    y=kx+b,
\]
где $x$ --- переменная, $k$ --- \textbf{угловой коэффициент},
$b$ --- \textbf{свободный коэффициент}.

\textbf{Что показывает $k$:}

\begin{itemize}
    \item $k>0$ --- функция возрастает:
    при движении слева направо прямая идёт вверх;

    \item $k<0$ --- функция убывает:
    при движении слева направо прямая идёт вниз;

    \item $k=0$ --- получаем горизонтальную прямую
    $y=b$, параллельную оси $Ox$.
\end{itemize}

\textbf{Что показывает $b$:}
число $b$ равно ординате точки пересечения графика
с осью $Oy$.

При $x=0$ всегда
\[
    y=b,
\]
поэтому график проходит через точку
\[
    (0;b).
\]

\begin{center}

\begin{tikzpicture}

\begin{axis}[
    width=0.78\linewidth,
    height=6.4cm,
    axis lines=middle,
    xmin=-4,
    xmax=5,
    ymin=-4,
    ymax=5,
    xtick={-3,-2,-1,1,2,3,4},
    ytick={-3,-2,-1,1,2,3,4},
    grid=both,
    minor tick num=0,
    xlabel={$x$},
    ylabel={$y$},
    clip=false,
    samples=200,
    unit vector ratio=1 1
]

\addplot[
    accent,
    very thick,
    domain=-4:5
]{0.7*x-1};

\addplot[
    darkaccent,
    dashed,
    very thick,
    domain=-4:5
]{-0.6*x+2};

\addplot[
    textgray,
    densely dotted,
    very thick,
    domain=-4:5
]{1.5};

\node[
    anchor=west,
    text=accent
]
at (axis cs:3.4,1.2)
{$k>0$};

\node[
    anchor=west,
    text=darkaccent
]
at (axis cs:2.8,0)
{$k<0$};

\node[
    anchor=west,
    text=textgray
]
at (axis cs:2.7,1.75)
{$k=0$};

\end{axis}

\end{tikzpicture}

\end{center}

\textbf{Правило чтения графика:}
направление возрастания или убывания определяем
\textbf{слева направо}.

% ================================================================
% 2
% ================================================================

\topic{2. Быстрое определение коэффициентов по графику}

Если на графике хорошо видны две точки
\[
    A(x_1;y_1)
    \qquad\text{и}\qquad
    B(x_2;y_2),
\]
то
\[
    k=
    \frac{\Delta y}{\Delta x}
    =
    \frac{y_2-y_1}{x_2-x_1}.
\]

То есть делим
\textbf{вертикальное изменение}
на
\textbf{горизонтальное}.

Знак коэффициента получается автоматически
из разностей координат.

\textbf{Важно:}
нельзя делить «горизонтальное на вертикальное».
Для углового коэффициента используется отношение
\[
    \frac{\text{изменение по }y}
         {\text{изменение по }x}.
\]

Свободный коэффициент $b$ удобно считывать
по точке пересечения графика с осью $Oy$.

Например, если эта точка имеет координаты
\[
    (0;-2),
\]
то
\[
    b=-2.
\]

\textbf{Пример.}

Пусть
\[
    A(1;3),
    \qquad
    B(6;1).
\]

Тогда
\[
    k=
    \frac{1-3}{6-1}
    =
    -\frac25.
\]

Подставим координаты точки $A$:
\[
    3=
    -\frac25\cdot1+b.
\]

Отсюда
\[
    b=
    3+\frac25
    =
    \frac{17}{5}.
\]

Следовательно,
\[
    \boxed{
        y=-\frac25x+\frac{17}{5}
    }.
\]

% ================================================================
% 3
% ================================================================

\topic{3. Метод узловых точек: точное восстановление формулы}

В рамках этого занятия
\textbf{узловой точкой}
будем называть точку графика,
обе координаты которой являются целыми числами.

Для функции
\[
    y=kx+b
\]
неизвестны два коэффициента:
$k$ и $b$.

Поэтому для восстановления формулы достаточно
взять \textbf{две} подходящие точки графика.

\textbf{Пример.}

Даны точки
\[
    A(2;1),
    \qquad
    B(6;2).
\]

Подставляем координаты в формулу
\[
    y=kx+b:
\]

\[
\begin{cases}
    1=2k+b,\\
    2=6k+b.
\end{cases}
\]

Вычитаем первое уравнение из второго:
\[
    1=4k,
\]
откуда
\[
    k=\frac14.
\]

Теперь найдём $b$:
\[
    b=
    1-2\cdot\frac14
    =
    \frac12.
\]

Искомая функция:
\[
    \boxed{
        y=\frac14x+\frac12
    }.
\]

\textbf{Самопроверка:}
после нахождения формулы подставьте
координаты обеих исходных точек.
Обе пары координат должны удовлетворять уравнению.

% ================================================================
% БЛОК-СХЕМА
% ================================================================

\clearpage

\thispagestyle{empty}

\begin{center}

{\Large\bfseries\color{darkaccent}
Алгоритм: найти $k$ и $b$ по двум точкам
}

\end{center}

\vspace{9mm}

\begin{center}

\begin{tikzpicture}[
    node distance=9mm and 16mm
]

\node[
    flowstart
]
(start)
{Начало};

\node[
    flowaction,
    below=of start
]
(pick)
{%
Выберите две точки графика\\
$A(x_1;y_1)$ и $B(x_2;y_2)$
};

\node[
    flowaction,
    below=of pick
]
(sub)
{%
Подставьте координаты в\\
$y=kx+b$ и составьте систему
};

\node[
    flowaction,
    below=of sub
]
(solve)
{%
Решите систему и найдите\\
$k$ и $b$
};

\node[
    flowaction,
    below=of solve
]
(write)
{%
Запишите формулу\\
$y=kx+b$
};

\node[
    flowdecision,
    below=13mm of write
]
(check)
{%
Обе точки удовлетворяют формуле?
};

\node[
    flowstart,
    below=13mm of check
]
(finish)
{Готово};

\node[
    flowaction,
    right=28mm of check
]
(error)
{%
Проверьте координаты,
знаки и арифметику
};

\draw[flowarrow]
(start)
--
(pick);

\draw[flowarrow]
(pick)
--
(sub);

\draw[flowarrow]
(sub)
--
(solve);

\draw[flowarrow]
(solve)
--
(write);

\draw[flowarrow]
(write)
--
(check);

\draw[flowarrow]
(check)
--
node[
    left,
    font=\small
]
{да}
(finish);

\draw[flowarrow]
(check)
--
node[
    above,
    font=\small
]
{нет}
(error);

\draw[flowarrow]
(error.north)
|-
([xshift=8mm]sub.east)
--
(sub.east);

\end{tikzpicture}

\end{center}

\clearpage

% ================================================================
% 4
% ================================================================

\topic{4. Семейство горизонтальных прямых $y=m$}

Формула
\[
    y=m
\]
задаёт семейство горизонтальных прямых,
параллельных оси $Ox$.

\begin{itemize}
    \item при $m=2$ получаем $y=2$;

    \item при $m=0$ получаем $y=0$,
    то есть ось $Ox$;

    \item при $m=-3$ получаем $y=-3$.
\end{itemize}

Изменение $m$ перемещает прямую
\textbf{вверх или вниз},
не меняя её наклона.

\begin{center}

\begin{tikzpicture}

\begin{axis}[
    width=0.76\linewidth,
    height=5.6cm,
    axis lines=middle,
    xmin=-4,
    xmax=4,
    ymin=-4,
    ymax=4,
    xlabel={$x$},
    ylabel={$y$},
    grid=both,
    samples=200,
    unit vector ratio=1 1,
    clip=false
]

\addplot[
    accent,
    very thick,
    domain=-4:4
]{2};

\addplot[
    darkaccent,
    very thick,
    domain=-4:4
]{0};

\addplot[
    textgray,
    very thick,
    domain=-4:4
]{-3};

\node[
    anchor=west,
    text=accent
]
at (axis cs:2.5,2.25)
{$y=2$};

\node[
    anchor=west,
    text=darkaccent
]
at (axis cs:2.5,0.25)
{$y=0$};

\node[
    anchor=west,
    text=textgray
]
at (axis cs:2.5,-2.75)
{$y=-3$};

\end{axis}

\end{tikzpicture}

\end{center}

% ================================================================
% 5
% ================================================================

\topic{5. Семейство прямых $y=kx$}

Если
\[
    b=0,
\]
то формула линейной функции принимает вид
\[
    y=kx.
\]

Все такие прямые проходят через начало координат
\[
    (0;0).
\]

Изменение $k$ меняет наклон прямой:

\begin{itemize}
    \item $k>0$ --- прямая возрастает;

    \item $k<0$ --- прямая убывает;

    \item $k=0$ --- получаем ось $Ox$;

    \item чем больше $|k|$,
    тем сильнее прямая отклоняется от оси $Ox$.
\end{itemize}

\begin{center}

\begin{tikzpicture}

\begin{axis}[
    width=0.76\linewidth,
    height=5.2cm,
    axis lines=middle,
    xmin=-3.5,
    xmax=3.5,
    ymin=-4.5,
    ymax=4.5,
    xlabel={$x$},
    ylabel={$y$},
    grid=both,
    samples=200,
    unit vector ratio=1 1,
    clip=false
]

\addplot[
    accent,
    very thick,
    domain=-3:3
]{x};

\addplot[
    darkaccent,
    very thick,
    domain=-2.2:2.2
]{2*x};

\addplot[
    textgray,
    dashed,
    very thick,
    domain=-1.5:1.5
]{-3*x};

\node[
    anchor=west,
    text=accent
]
at (axis cs:2.2,2.3)
{$k=1$};

\node[
    anchor=west,
    text=darkaccent
]
at (axis cs:1.5,3.2)
{$k=2$};

\node[
    anchor=west,
    text=textgray
]
at (axis cs:-1.3,3.9)
{$k=-3$};

\end{axis}

\end{tikzpicture}

\end{center}

% ================================================================
% 6
% ================================================================

\topic{6. Линейная функция в задачах по физике}

Формулы движения имеют ту же структуру,
что и
\[
    y=kx+b.
\]

\textbf{Скорость при равноускоренном движении:}
\[
    v(t)=v_0+at.
\]

Здесь свободный коэффициент --- $v_0$,
а угловой коэффициент --- ускорение $a$.

На графике $v(t)$:
\[
    v_0=v(0),
    \qquad
    a=\frac{\Delta v}{\Delta t}.
\]

\textbf{Координата при равномерном движении:}
\[
    x(t)=x_0+vt.
\]

Здесь свободный коэффициент ---
начальная координата $x_0$,
а угловой коэффициент --- скорость $v$:

\[
    x_0=x(0),
    \qquad
    v=\frac{\Delta x}{\Delta t}.
\]

Например, если за $4$ с координата увеличилась
на $2$ м, то
\[
    v=
    \frac{2}{4}
    =
    0{,}5\ \text{м/с}.
\]

% ================================================================
% 7
% ================================================================

\clearpage

\topic{7. Условие пересечения графиков}

В точке пересечения двух графиков
значения $x$ и $y$ у них совпадают.

Поэтому для функций
\[
    y=k_1x+b_1
\]
и
\[
    y=k_2x+b_2
\]
координату точки пересечения по $x$
находят из равенства
\[
    k_1x+b_1
    =
    k_2x+b_2.
\]

После нахождения $x$
его подставляют в любую из двух формул
и получают соответствующее значение $y$.

Для двух прямых полезно помнить:

\begin{itemize}
    \item если
    \[
        k_1\ne k_2,
    \]
    прямые пересекаются в одной точке;

    \item если
    \[
        k_1=k_2,
        \qquad
        b_1\ne b_2,
    \]
    прямые параллельны;

    \item если
    \[
        k_1=k_2,
        \qquad
        b_1=b_2,
    \]
    графики совпадают.
\end{itemize}

% ================================================================
% ТРЕНИРОВОЧНЫЙ БЛОК
% ================================================================

\clearpage

\topic{Тренировочный блок}

\textbf{Путь А.}

Выполните 5 заданий.
Решения оформляйте последовательно.
В задачах на восстановление формулы
обязательно записывайте итоговый вид функции.

\begin{enumerate}

    \item
    Для функции
    \[
        y=-3x+4
    \]
    укажите:

    а) возрастает она или убывает;

    б) координаты точки пересечения
    с осью $Oy$.

    \item
    Прямая проходит через точки
    \[
        A(1;2)
        \qquad\text{и}\qquad
        B(5;6).
    \]
    Найдите угловой коэффициент $k$.

    \item
    По двум точкам
    \[
        A(2;3)
        \qquad\text{и}\qquad
        B(6;1)
    \]
    восстановите формулу линейной функции
    \[
        y=kx+b.
    \]

    \item
    Для равномерного движения координата тела
    задаётся линейной зависимостью от времени.

    Известно, что график проходит через точки
    \[
        (0;-3)
        \qquad\text{и}\qquad
        (8;1),
    \]
    где время измеряется в секундах,
    а координата --- в метрах.

    Найдите начальную координату $x_0$
    и скорость $v$.

    \item
    Найдите точку пересечения графиков
    \[
        y=2x-5
    \]
    и
    \[
        y=-x+4.
    \]

\end{enumerate}

% ================================================================
% ОТВЕТЫ
% ================================================================

\clearpage

\topic{Ответы}

\renewcommand{\arraystretch}{1.35}

\begin{table}[h]

\centering

\begin{tabularx}{\linewidth}{
    @{}
    >{\bfseries}c
    X
    @{}
}

\toprule

№ & Ответ \\

\midrule

1 &
а) убывает;
б) $(0;4)$.
\\

2 &
$k=1$.
\\

3 &
\[
    k=-\frac12,
    \qquad
    b=4,
\]
поэтому
\[
    y=-\frac12x+4.
\]
\\

4 &
\[
    x_0=-3\ \text{м},
\]
\[
    v=
    \frac{1-(-3)}{8}
    =
    0{,}5\ \text{м/с}.
\]
\\

5 &
$(3;1)$.
\\

\bottomrule

\end{tabularx}

\end{table}

\end{document}